\documentclass[11pt,a4paper]{article}
\usepackage[margin=1in]{geometry}
\usepackage{booktabs}
\usepackage{longtable}
\usepackage{array}
\usepackage{amsmath}
\usepackage{listings}
\usepackage{xcolor}
\usepackage{hyperref}

\hypersetup{colorlinks=true, linkcolor=black, urlcolor=blue, citecolor=black}

\lstset{
  basicstyle=\ttfamily\small,
  keywordstyle=\color{blue!60!black},
  commentstyle=\color{green!40!black},
  stringstyle=\color{orange!50!black},
  showstringspaces=false,
  breaklines=true,
  frame=single
}

\title{Project 3 Report\\\large Matrix Multiplication in C and Performance Optimization}
\author{}
\date{}

\begin{document}
\maketitle

\section{Project Topic}
This project implements dense matrix multiplication in C and compares several optimization strategies on square matrices of different sizes.

\section{Experimental Environment}
\begin{itemize}
    \item Operating system: Linux / WSL2
    \item Kernel: \texttt{Linux ztw 6.6.87.2-microsoft-standard-WSL2 x86\_64}
    \item Compiler: \texttt{cc (Ubuntu 13.3.0-6ubuntu2\textasciitilde24.04.1) 13.3.0}
    \item Compiler flags: \texttt{-std=c11 -O3 -Wall -Wextra -pedantic}
    \item Parallel support: OpenMP
    \item BLAS library: OpenBLAS
\end{itemize}

The benchmark in this report was built and executed with:

\begin{lstlisting}[language=bash]
make clean
make USE_OPENBLAS=1 USE_OPENMP=1
./benchmark > 1.out
\end{lstlisting}

\section{Implemented Methods}
The project includes five multiplication paths:
\begin{itemize}
    \item \texttt{plain}: baseline triple-loop implementation
    \item \texttt{v1}: loop reordering from \texttt{ijk} to \texttt{ikj}
    \item \texttt{v2}: blocked multiplication with OpenMP on top of \texttt{ikj}
    \item \texttt{v3}: Strassen recursive multiplication
    \item \texttt{openblas}: \texttt{cblas\_sgemm} as the external reference
\end{itemize}
All matrix elements use \texttt{float}.

\section{Baseline Implementation: plain}
The baseline version uses the standard triple nested loop:

\begin{lstlisting}[language=C]
for (i = 0; i < a->rows; ++i)
    for (j = 0; j < b->cols; ++j)
        for (k = 0; k < a->cols; ++k)
            c->data[i * c->cols + j] +=
                a->data[i * a->cols + k] * b->data[k * b->cols + j];
\end{lstlisting}

This version is simple and easy to explain, so it is a good baseline. Its main drawback is poor cache locality when accessing \texttt{B[k][j]} by column.

\section{Improved Implementations}
\subsection{v1: Loop Reordering}
The idea of \texttt{v1} is to change the loop order from \texttt{ijk} to \texttt{ikj}. This allows the innermost loop to access one full row of \texttt{B} and one full row of \texttt{C} consecutively.

\begin{lstlisting}[language=C]
for (i = 0; i < a->rows; ++i) {
    const float *a_row = &a->data[i * a->cols];
    float *c_row = &c->data[i * c->cols];

    for (k = 0; k < a->cols; ++k) {
        const float a_ik = a_row[k];
        const float *b_row = &b->data[k * b->cols];

        for (j = 0; j < b->cols; ++j)
            c_row[j] += a_ik * b_row[j];
    }
}
\end{lstlisting}

Main effects of \texttt{v1}:
\begin{itemize}
    \item more contiguous memory access for \texttt{B} and \texttt{C}
    \item fewer cache misses than the baseline
    \item more friendly to compiler auto-vectorization
\end{itemize}

This version is not a hand-written SIMD intrinsic implementation. However, it is SIMD-friendly. GCC vectorization reports show that the innermost \texttt{j} loop in \texttt{v1} was automatically vectorized under \texttt{-O3 -fopenmp}:

\begin{lstlisting}
matmul.c:91:27: optimized: loop vectorized using 16 byte vectors
matmul.c:91:27: optimized: loop versioned for vectorization because of possible aliasing
\end{lstlisting}

\subsection{v2: Blocking + OpenMP}
\texttt{v2} adds cache blocking on top of the \texttt{ikj} traversal. The matrix is split into smaller tiles with block size \texttt{64}.

\begin{lstlisting}[language=C]
for (i0 = 0; i0 < a->rows; i0 += MATMUL_BLOCK_SIZE) {
    size_t i1 = i0 + MATMUL_BLOCK_SIZE < a->rows ?
                i0 + MATMUL_BLOCK_SIZE : a->rows;

    for (k0 = 0; k0 < a->cols; k0 += MATMUL_BLOCK_SIZE) {
        size_t k1 = k0 + MATMUL_BLOCK_SIZE < a->cols ?
                    k0 + MATMUL_BLOCK_SIZE : a->cols;

        for (j0 = 0; j0 < b->cols; j0 += MATMUL_BLOCK_SIZE) {
            size_t j1 = j0 + MATMUL_BLOCK_SIZE < b->cols ?
                        j0 + MATMUL_BLOCK_SIZE : b->cols;

            for (i = i0; i < i1; ++i)
                for (k = k0; k < k1; ++k)
                    for (j = j0; j < j1; ++j)
                        c_row[j] += a_ik * b_row[j];
        }
    }
}
\end{lstlisting}

OpenMP is applied to the outer block loop:

\begin{lstlisting}[language=C]
#pragma omp parallel for schedule(static) private(j0, k0)
\end{lstlisting}

Main effects of \texttt{v2}:
\begin{itemize}
    \item improved cache locality through blocking
    \item parallel execution through OpenMP
    \item much better performance on medium and large matrices
\end{itemize}

GCC vectorization reports also show auto-vectorization inside the block kernel:

\begin{lstlisting}
matmul.c:134:40: optimized: loop vectorized using 16 byte vectors
matmul.c:134:40: optimized: loop versioned for vectorization because of possible aliasing
\end{lstlisting}

So the speedup of \texttt{v2} comes from blocking, OpenMP, and compiler auto-vectorization.

\subsection{v3: Strassen Recursive Algorithm}
\texttt{v3} uses the Strassen algorithm, which is an algorithm-level optimization rather than a simple loop optimization. A matrix is split into four submatrices:

\[
A = \begin{bmatrix} A_{11} & A_{12} \\ A_{21} & A_{22} \end{bmatrix},
\quad
B = \begin{bmatrix} B_{11} & B_{12} \\ B_{21} & B_{22} \end{bmatrix}
\]

Strassen computes seven intermediate products:
\[
\begin{aligned}
M_1 &= (A_{11} + A_{22})(B_{11} + B_{22}) \\
M_2 &= (A_{21} + A_{22})B_{11} \\
M_3 &= A_{11}(B_{12} - B_{22}) \\
M_4 &= A_{22}(B_{21} - B_{11}) \\
M_5 &= (A_{11} + A_{12})B_{22} \\
M_6 &= (A_{21} - A_{11})(B_{11} + B_{12}) \\
M_7 &= (A_{12} - A_{22})(B_{21} + B_{22})
\end{aligned}
\]

Then the result blocks are reconstructed as:
\[
\begin{aligned}
C_{11} &= M_1 + M_4 - M_5 + M_7 \\
C_{12} &= M_3 + M_5 \\
C_{21} &= M_2 + M_4 \\
C_{22} &= M_1 - M_2 + M_3 + M_6
\end{aligned}
\]

The implementation follows this structure directly. A representative code fragment is:

\begin{lstlisting}[language=C]
if ((rc = matrix_add(a11, a22, s1, 1)) != 0) goto cleanup;
if ((rc = matrix_add(b11, b22, s2, 1)) != 0) goto cleanup;
if ((rc = matmul_strassen_recursive(s1, s2, m1)) != 0) goto cleanup;
...
if ((rc = matrix_add(m1, m4, c11, 1)) != 0) goto cleanup;
if ((rc = matrix_add(c11, m5, c11, -1)) != 0) goto cleanup;
if ((rc = matrix_add(c11, m7, c11, 1)) != 0) goto cleanup;
\end{lstlisting}

Its theoretical complexity is lower than classical \(O(n^3)\), but the actual implementation introduces heavy recursion, temporary matrix allocation, and data movement. Therefore, it does not necessarily outperform \texttt{v2} in practice.

\section{Correctness Check}
According to \texttt{1.out}, all small validation cases passed:
\begin{itemize}
    \item \texttt{plain vs v1}: \texttt{n=4,7,16,31} all PASS
    \item \texttt{plain vs v2}: \texttt{n=4,7,16,31} all PASS
    \item \texttt{plain vs v3}: \texttt{n=4,7,16,31} all PASS
    \item \texttt{v3 vs OpenBLAS}: \texttt{n=4,7,16,31} all PASS
\end{itemize}

This shows that all three improved versions produced correct results in the tested cases, and \texttt{v3} matched OpenBLAS as well.

\section{Benchmark Results}
\subsection{Results for \(n = 16\)}
\begin{center}
\begin{tabular}{lrrl}
\toprule
Implementation & Time (s) & GFLOPS & Correctness \\
\midrule
plain    & 0.000003    & 2.990    & PASS \\
v1       & 0.000001    & 6.776    & PASS \\
v2       & 0.000042    & 0.194    & PASS \\
v3       & 0.000042    & 0.194    & PASS \\
openblas & 0.000000    & 16.483   & PASS \\
\bottomrule
\end{tabular}
\end{center}

\subsection{Results for \(n = 128\)}
\begin{center}
\begin{tabular}{lrrl}
\toprule
Implementation & Time (s) & GFLOPS & Correctness \\
\midrule
plain    & 0.001897    & 2.211    & PASS \\
v1       & 0.000326    & 12.856   & PASS \\
v2       & 0.000645    & 6.502    & PASS \\
v3       & 0.000246    & 17.032   & PASS \\
openblas & 0.000038    & 110.009  & PASS \\
\bottomrule
\end{tabular}
\end{center}

\subsection{Results for \(n = 1024\)}
\begin{center}
\begin{tabular}{lrrl}
\toprule
Implementation & Time (s) & GFLOPS & Correctness \\
\midrule
plain    & 2.227743    & 0.964     & PASS \\
v1       & 0.087690    & 24.490    & PASS \\
v2       & 0.019736    & 108.811   & PASS \\
v3       & 0.291508    & 7.367     & PASS \\
openblas & 0.002139    & 1004.035  & PASS \\
\bottomrule
\end{tabular}
\end{center}

\subsection{Results for \(n = 8192\)}
\begin{center}
\begin{tabular}{lrrl}
\toprule
Implementation & Time (s) & GFLOPS & Correctness \\
\midrule
plain    & 7013.361078  & 0.157    & PASS \\
v1       & 113.915864   & 9.652    & PASS \\
v2       & 8.246949     & 133.323  & PASS \\
v3       & 66.821979    & 16.454   & PASS \\
openblas & 1.632326     & 673.586  & PASS \\
\bottomrule
\end{tabular}
\end{center}

\subsection{Comparison of the Three Improved Versions}
\begin{center}
\begin{tabular}{rrrrrrr}
\toprule
Size & v1 Time (s) & v1 GFLOPS & v2 Time (s) & v2 GFLOPS & v3 Time (s) & v3 GFLOPS \\
\midrule
16   & 0.000001   & 6.776   & 0.000042   & 0.194    & 0.000042   & 0.194 \\
128  & 0.000326   & 12.856  & 0.000645   & 6.502    & 0.000246   & 17.032 \\
1024 & 0.087690   & 24.490  & 0.019736   & 108.811  & 0.291508   & 7.367 \\
8192 & 113.915864 & 9.652   & 8.246949   & 133.323  & 66.821979  & 16.454 \\
\bottomrule
\end{tabular}
\end{center}

\subsection{Result for \(n = 65536\)}
\begin{center}
\begin{tabular}{ll}
\toprule
Implementation & Result \\
\midrule
All implementations & allocation failed \\
\bottomrule
\end{tabular}
\end{center}

The output reports an estimated working set of \texttt{65536.00 MiB}, which could not be allocated in the current environment.

\section{Result Analysis}
\begin{enumerate}
    \item The baseline \texttt{plain} version is always the slowest one. Its performance is especially poor for \(1024\) and \(8192\), which confirms that the original triple loop has poor cache locality on large matrices.
    \item \texttt{v1} is clearly faster than \texttt{plain}, showing that loop reordering alone already improves memory access behavior.
    \item \texttt{v2} is the best self-implemented version in this project. It reaches \texttt{108.811 GFLOPS} at \(1024\) and \texttt{133.323 GFLOPS} at \(8192\), much faster than \texttt{plain}, \texttt{v1}, and \texttt{v3}.
    \item \texttt{v3} uses the Strassen recursive algorithm. Although its theoretical complexity is better, the current implementation suffers from recursive overhead, temporary matrix allocation, and extra matrix addition/subtraction. As a result, it is slower than \texttt{v2} in practice.
    \item The comparison table of \texttt{v1}, \texttt{v2}, and \texttt{v3} shows that small-size results fluctuate more, but for representative sizes such as \(1024\) and \(8192\), \texttt{v2} is clearly the best among the three self-implemented optimized versions.
    \item OpenBLAS is faster than every self-implemented version at all meaningful sizes, which is expected because it uses mature packing, vectorization, and threading strategies.
\end{enumerate}

\section{Conclusion}
This project implemented four matrix multiplication methods: \texttt{plain}, \texttt{v1}, \texttt{v2}, and \texttt{v3}, and compared them with OpenBLAS.

The benchmark results show:
\begin{itemize}
    \item \texttt{v1} confirms that loop reordering is an effective optimization.
    \item \texttt{v2} demonstrates that blocking is the most successful optimization in this project for medium and large matrices.
    \item \texttt{v3}, although based on the Strassen algorithm, does not outperform \texttt{v2} in the current implementation.
    \item OpenBLAS remains the fastest overall implementation.
\end{itemize}

Therefore, if the criterion is the best self-implemented method, the best result in this project is \texttt{v2}. If the criterion is overall peak performance, OpenBLAS is the best performer.

\end{document}
